\documentclass[10pt]{article}

\usepackage[utf8]{inputenc}

\usepackage[T1]{fontenc}

\usepackage{amsmath}

\usepackage{amsfonts}

\usepackage{amssymb}

\usepackage[version=4]{mhchem}

\usepackage{stmaryrd}

\usepackage{bbold}



\begin{document}

существует $V \in Y$, притом только один такой, что $(u, v) \subset M$, тогда $f_{M}(u)=v$.



Два отображения $f$ и $g$ наз. р а в ны м и, если области их определения совпадают и $f(x)=g(x)$ для любого $x \in X$. В этом случае совпадают и области значений этих О. Отображение $f$ на $X$ наз. по о т о я н н ы м, если $f(x)=a$ для любого $x \in X$. С у ж е н и е м о т о б$\mathbf{p}$ а же в и я $f: X \rightarrow Y$ на подмножество $A \subset X$ наз. отображение $q$, заданное на множестве $A$ равенством $\varphi(x)=f(x), x \in A$; это сужение обозначается $f_{A}$. Р а cшш и $\mathrm{p}$ ен и е м о тоб р а ения (или продолже н и е м) на множество $E \supset X$ наз. отображение $F$, определенное на $E$ и удовлетворяющее равенству $F(x)=f(x)$ для всех $x \in X$. Если заданы три множества $X, Y, Z$, на $X$ определено отображение $f$ со значениями в $Y$, а на $Y$ задано отображение $g$ со значениями в $Z$, то существует отображение $h$ с областью определения $X$, принимающее значения в $Z$ и определяемое равенством $h(x)=g[f(x)]$. Это О. наз. к о м п о з и ц й $е$ й отображений $f$ и $g$, a $f$ и $g$ соответствую щ и м и ото об р жения м и, и обозначается $g \circ f$, причем порядок записи играет существенную роль (для функций действительного переменного принят термин с у п е р по з и ц и я). Отображение $h$ наз. также с ло жны м о то б р а же н и е м, составленным из внутреннего отображения $f$ и внеппего отображения $g$. Понятие сложного О. обобщается на любое конечное число составляющих 0 .



Отображение $f$, определенное на $X$ и принимающее значения в $Y$, порождает новое $0 .$, заданное на подмножествах множества $X$ и пмеющее в качестве значений подмножества множества $Y$. Именно, если $A \subset X$, то



\[

f(A) \stackrel{\text { def }}{=}\{y=f(x), x \in A\} .

\]



Множество $f(A)$ наз. о б р а з о м м н о же с т в а $A$. Если положить $A=\{x\}$, то получится исходное отображение $f(x)$, так что $f(A)$ есть расширение отображения $f(x)$ с множества $X$ на множество $\mathfrak{P}(X)$ всех подмножеств множества $X$, если отождествлять одноәлементное множество с элементом, его составляющим. В случае $Y=X$ множество $A$ наз. и н в а р и а н тны м по д м н о же с т в о м отображения $f$, если $f(A) \subset A$, а точка $x$ наз. н е по о д в и ж н ы м ә л ем е в т о м отображения $f$, если $f(x)=x$. Иввариантные множества и неподвижные элементы играют важвую роль при решении функциональных уравнений вида $f(x)=\boldsymbol{a}$ или $x-f(x)=\boldsymbol{a}$.



Каждое отображение $f: X \rightarrow Y$ порождает О., заданное на подмножествах множества $f(X)$ и имеющее в качестве значений подмножества множества $X$. Именно, для каждого $B \subset f(X)$ через $f^{-1}(B)$ обозначается множество $\{x, f(x) \in B\}$, называемое по л ны м п р оо б р а з м м множества $B$. Если $f^{-1}(y)$ для любого $y \in f(X)$ состоит из единственного элемента, то $f^{-1}$ есть $О$. элемента, определенное на $f(X)$, принимающее значения в $X$ и называемое о $б$ р а т н ы м о то б р а ж е в и е м к $f$. Существование обратного 0 . эквивалентно разрешимости уравнения $f(x)=y$, $y \in f(X)$.



Если множества $X$ и $Y$ наделены нек-рыми свойствами, то во множестве $F(X, Y)$ всех 0 . на $X$ в $Y$ могут быть выделены содержательные классы. Так, для частично упорядоченных множеств $X$ и $Y$ отображение $f$ наз. и зо то н ны м, если $x<y$ влечет $f(x) \leqslant f(y)$. Для комплексных плоскостей $X$ и $Y$ выделяется класс голоморфых 0 . Для топологич. пространств $X$ и $Y$ естественным образом выделяется класс непрерывных $O$. этих пространств; строится развернутая теория дифференцирования отображений. Для О. скалярного аргумента и, в более общем случае, для О., определенных на пространстве с мерой, может быть введено по- нятие (сильной или слабой) пзмеримости и могут быть построены различные интегралы лебеговского типа (напр., Бохнера интеграл, Даниеля интеграл).



О. наз. многознатным о тобра жен и ем, если нек-рым значениям $x$ соответствуот подмножества $\boldsymbol{Y}_{x} \subset Y$, состояиие более чем из одного элемента. Таковы, напр., многолистные функции комплексного переменного, многозначные О. топологических пространств й др.



Лum.: [1] Б у р б ак и Н., Теория множеств, пер. с франц., М., 1965; [2] е г о ж е, Общая топология. Основные структуры, пер. с франц., 2 изд., М., 1968; [3] К е л ли Д ж. Л., Общая то-

пология, пер. с англ., 2 изд., М., 1981.

В. И. Соболев.



ОТОБРАЖЕНИЕ ПЕРИОДОВ - отображение, сопоставляющее точке $s$ базы $S$ семейства $\left\{X_{s}\right\}$ алгебраич. многообразий над полем $\mathbb{C}$ комплексных чисел когомологии $H^{*}\left(X_{s}\right)$ слоя над этой точкой, снабженные $Х о д ж а$ структурой. Полученная при этом структура Ходжа рассматривается как точка в многообразии модулей структур Ходжа данного типа.



Изучение О. п. восходит к исследованиям Н. Абеля (N. Abel) и К. Якоби (C. Jacobi) интегралов алгебракч. функций (см. Абелев дифференциал). Однако до недавнего времени глубоко были изучены лишь О. П., отвечающие семействам кривых.



Пусть $\left\{X_{s}\right\}$ есть семейство слоев $X_{s}=f^{-1}(s)$ гладкого проективного морфизма $f: X \rightarrow S$, где $S$ - гладкое многообразие. Тогда когомологии $H^{*}\left(X_{s}, \mathbb{Z}\right)=V_{\mathbb{Z}}$ снабжены чистой, поляризованной структурой Ходжа, к-рая задается гомоморфизмом вещественных аләебраических әруnn $h: \mathbb{C}^{*} \rightarrow G_{\mathbb{R}}$, где $\mathbb{C}^{*}-$ мультипликативная группа $\mathbb{C}^{*}$ поля комплексных чисел, рассматриваемая как вещественная алгебраич. группа, а



\[

G=\{g \in \mathrm{GL}(V), \psi(g x, g y)=\lambda(g) \psi(x, y)\}

\]



\begin{itemize}

  \item алгебраич. группа линейных преобразований пространства $V$, умножающих невырожденную (симметрическую или кососимметрическую) билинейную форму $\psi$ на скалярный множитель; причем автоморфизм Ad $h(i)$ группы $G_{\mathbb{R}}$ является инволюциейі Картана и $h(\mathbb{R} *$ лежит в центре группы $G_{R}$. Множество $X_{G}$ гомоморфизмов $h: \mathbb{C}^{*} \rightarrow G_{\mathbb{R}}$, обладающих указанными свойствами, естественным образом снабжено $G_{R}$-инвариантной структурой однородного кәлерова многообразия и наз. м ног ооб разие м $\Gamma$ и иффи т с а, а фактор $M_{G}=X_{G} / G_{\mathbb{Z}}$ является пространством модулей структур Ходжа. Гомоморфизм $h$ задает разложение Ходжа

\end{itemize}



\[

\text { (S) }_{\mathbb{C}}=\bigoplus \text { (S) }^{p,-p}

\]



алгебры Ли (S) группы $G$, где (\$) $p,-p-$ подпространство в (\$) $\mathbb{C}$, на к-ром Ad $h(z)$ действует умножением на $\bar{z} p_{z}^{-p}$. Сопоставление $h \rightarrow P(h)$, где $P(h)-$ параболич. подгруппа в $G_{\mathbb{C}}$, алгебра Ли к-рой есть $\oplus_{p \geqslant 0} \mathfrak{S}^{p,-p}$, задает открытое плотное вложение многообразия $X_{G}$ в компактное $G_{\mathbb{C}}$-однородное многообразие флагов $\check{X}_{G}$. В касательном пространстве



\[

\mathfrak{G}_{\mathbb{G}} / \oplus_{p \geqslant 0} \mathbb{S}^{p,-p}

\]



к $X_{G}$ в точке $h$ выделено $\mathbf{\Gamma} 0 \mathrm{p}$ и 3 о в т а л в н о е по о д п р о с т р а н с т в о



\[

\oplus_{p \geqslant-1} \mathbb{S}^{p,-p} \mid \bigoplus_{p \geqslant 0 \mathbb{S}^{p},-p}

\]



Голоморфное отображение в $X_{G}$ или $M_{G}$ наз. г о р из о н т а л ь в ы м, если образ его касательного отпбражения лежит в горизонтальном подрасслоении.





\end{document}